% ---------------- 原书前言（Preface，牛津 2001） ----------------
\chapter*{前言}
\addcontentsline{toc}{chapter}{前言}

\begin{quote}
\centering \kaishu “……万有也靠他而立。”\quad（《歌罗西书》一章十七节）
\end{quote}

磁学是一门被研究了近三千年的学问。磁石（lodestone）——一种铁矿石——最早引起了希腊
学者与哲学家的注意，而导航用磁性罗盘则是这项研究结出的第一件技术成果。罗盘固然至迟在公元十二世纪已为西欧所知，但直到约 1600 年，才出现堪称近代水平的对罗盘原理的
解释。最近两个世纪以来，进展大大加快，并形成了两项把磁学与其他物理现象联系起来的
重大成果。其一，磁与电密不可分，二者共同构成了光——所谓电磁波。其二，这一联系的
根源在于相对论，因此磁学可以说是一种纯粹的相对论性效应，它源于观察者与导线中、
乃至铁原子中运动电荷的相对运动。然而，时至今日真正引人入胜的，仍是凝聚态体系——
包括铁磁体、自旋玻璃与低维体系——中的磁性。宏观体系所表现出的磁性质，与原子、
分子有着根本的不同，尽管二者由完全相同的基本组分构成。其原因在于，磁是一种集体
现象，需要极大量粒子的彼此协同；在这个意义上，它与超导、超流、乃至固态现象本身
颇为相似。对这类基本问题的探究，与寻找新型永磁材料、传感器与记录应用材料的技术
驱动，始终并行不悖。

本书源于作者在牛津大学为选择了凝聚态物理方向的三年级、四年级本科生开设的一门
课程。当时显而易见，我们需要一本既讲述基础、又提供背景材料与课堂上无法涵盖的
补充专题的教材。本书的目标，是把这一领域作为一个连贯的整体呈现出来，提供有用而
有趣的原始素材，并且读来令人愉悦。本书同时也是《牛津凝聚态物理 master 系列》
（Oxford Master Series in Condensed Matter Physics）的一种；该系列其他各卷分别
讨论电子性质、光学性质、超导电性、结构与软凝聚态物质。

阅读本书的前提是掌握基本的量子力学与电磁学，并熟悉原子物理的若干结果。
这些内容已总结于书末附录，既便于读者查阅，也统一了全书的记号。

凝聚态体系中那些有趣的磁效应有两个关键要素：其一，原子必须具有磁矩；其二，
这些磁矩之间必须以某种方式发生相互作用。这两个主题分别在第二章与第四章中讨论。
第二章回答“原子为什么具有磁矩？”这一问题，并说明磁矩彼此并无相互作用时它们的行为
及研究方法。\mnote{本课程的一些可能安排：（1）短课程（假定已知第 1 章内容）：
第 2 章（略去 2.6--2.8）；第 3 章（略去 3.2）；第 4 章（略去 4.2.5、4.2.6）；
第 5 章（略去 5.4--5.7）；第 6 章（略去 6.4--6.5）；第 7 章（略去 7.1--7.2）。
（2）较长课程：第 1--4 章；第 5 章（5.6、5.7 作背景阅读）；第 6 章；第 7 章
（7.5--7.9 作背景阅读）；第 8 章选讲专题。}第三章描述晶体内的局域环境如何影响这些磁矩，以及研究这种影响的各种
技术。第四章则回答“不同原子上的磁矩如何彼此相互作用？”。有了这两个要素，
磁有序便可能出现，这是第五、六章的主题。第五章描述固体中出现的各种类型的磁有序。
第六章再次讨论有序问题，但从对称性破缺的基本观念出发，描述相变、激发与磁畴。
本章特别强调磁有序与超导电性等其他类型的对称性破缺基态之间的联系。第七章专门
讨论金属的磁性质——在金属中，磁性往往与巡游的传导电子相关。第八章描述当存在
相互竞争的磁相互作用、或体系维度降低时出现的一些精微而复杂效应。这些专题至今
仍是研究热点，尚有许多悬而未决的问题。全书始终贯穿对材料性质与应用的讨论，
以展示这些思想对真实材料的意义——包括铁氧体、永磁体，以及近年来具有巨大技术
价值的各种磁光效应与磁电阻效应背后的物理。这是一本为物理学家而写的书，因此
重点放在支撑整个领域的清晰的物理原理上：量子力学、对称性与电磁学。但它并不只是
一本“理论书”，而是力求把这一学科与实验物理学家当前实际使用的测量手段与实验
技术联系起来，在本科基础物理的原理与当前研究热点之间架起一座桥梁。

第一章至第七章的末尾附有延伸阅读书目与习题。习题难度不一，用以深化正文中
讨论的问题。第八章未设习题（该章所述皆为当前研究课题），但提供了详尽的延伸
阅读书目。

本书写作过程中得到了许多人的帮助，在此谨致谢忱。感谢牛津大学出版社的
Sönke Adlung 及其团队的支持，也感谢本 master 系列的其他作者。牛津大学
Mansfield 学院与牛津大学物理系为我提供了令人振奋的工作环境。我要感谢我的
学生们——他们常常促使我对自己原以为已经理解的东西作更深入的思考。在本书
各个方面的准备过程中，我与 Hideo Aoki、Arzhang Ardavan、Deepto Chakrabarty、
Amalia Coldea、Radu Coldea、Roger Cowley、Steve Cox、Gillian Gehring、
Matthias Gester、John Gregg、Martin Greven、Mohamedally Kurmoo、Steve Lee、
Wilson Poon、Francis Pratt、John Singleton 与 Candadi Sukumar 的讨论使我
受益良多。我尤其要感谢那些阅读了本书各轮草稿的挚友与同行，他们严苛的批评与
富有洞见的问题使最终成品大为改进：Katherine Blundell、Richard Blundell、
Andrew Boothroyd、Geoffrey Brooker、Bill Hayes、Brendon Lovett、
Lesley Parry-Jones 与 Peter Riedi。本书付印后若发现错误，将公布于本书网页：
\begin{center}
\texttt{http://users.ox.ac.uk/\textasciitilde sjb/magnetism/}
\end{center}

最重要的是，我要感谢 Katherine——我亲爱的妻子与灵魂伴侣。她给予我的启迪、
忠告、友谊与爱，远胜他人。本书谨献给她。

\begin{flushright}
牛津\quad 2001 年 5 月\\
S.~J.~B.
\end{flushright}

\clearpage
